%----------------------------------------------------------------------------------------
%	COVER LETTER --- Goldman Sachs Summer Analyst (Engineering, 2027 batch)
%	Companion to resume.tex
%----------------------------------------------------------------------------------------
\documentclass[a4paper,11pt]{article}
\usepackage{url}
\usepackage{parskip}
\RequirePackage{color}
\RequirePackage{graphicx}
\usepackage[usenames,dvipsnames]{xcolor}
\usepackage[top=1.5cm, bottom=1.5cm, left=1.8cm, right=1.8cm]{geometry}
\usepackage{tabularx}
\newcolumntype{C}{>{\centering\arraybackslash}X}
\usepackage[unicode, draft=false]{hyperref}
\definecolor{linkcolour}{rgb}{0,0.2,0.6}
\hypersetup{colorlinks,breaklinks,urlcolor=linkcolour,linkcolor=linkcolour}
\usepackage{fontawesome5}
\setlength{\parskip}{6pt}

\begin{document}
\pagestyle{empty}

%----------------------------------------------------------------------------------------
%	HEADER
%----------------------------------------------------------------------------------------
\begin{tabularx}{\linewidth}{@{} C @{}}
{\Large \textbf{Manuvarthi Seshadri Naidu}} \\[4pt]
\href{https://sesh-s-portfolio.vercel.app}{\raisebox{-0.05\height}\faGlobe\ sesh-s-portfolio.vercel.app} \ $|$ \
\href{https://github.com/Seshmanuvarthi}{\raisebox{-0.05\height}\faGithub\ Seshmanuvarthi} \ $|$ \
\href{https://www.linkedin.com/in/seshadri-naidu-manuvarthi-366466295/}{\raisebox{-0.05\height}\faLinkedin\ seshadri-naidu-manuvarthi} \ $|$ \
\href{mailto:seshmanuvarthi27@gmail.com}{\raisebox{-0.05\height}\faEnvelope\ seshmanuvarthi27@gmail.com} \ $|$ \
\href{tel:+918897055099}{\raisebox{-0.05\height}\faMobile\ +91 8897055099} \\
\end{tabularx}

\vspace{14pt}

\noindent May 21, 2026

\vspace{4pt}

\noindent Hiring Team \\
Goldman Sachs --- Engineering Division \\
Bengaluru / Hyderabad, India

\vspace{6pt}

\noindent \textbf{Re: Summer Analyst, Engineering (2027 batch, Bengaluru / Hyderabad)}

\vspace{8pt}

Dear Hiring Team,

I'm applying for the 2027 Summer Analyst program in Goldman Sachs Engineering. One line from your description stuck with me --- \textit{``Our Engineers don't just make things, we make things possible.''} That framing matches how I've come to think about my own work: the interesting part isn't the demo, it's the engineering decisions that decide whether a system holds up when something goes wrong.

The clearest example is the SSH honeypot I built last semester. It captures live attacker sessions into a normalised SQLite schema and routes unknown commands to an LLM-powered classifier (Groq API, Llama-3.3-70B-versatile). The interesting design work wasn't about the model --- it was about making the system trustworthy when the model could be wrong. I added structured-JSON output enforcement, a 3-tier parse-and-repair pipeline, and confidence-gated acceptance ($\geq$0.75) so the rule engine could fall back to the model only when it was confident, and otherwise return a deterministic ``unknown'' rather than a hallucinated label. A chaos-engine layer then escalates response intensity on repeated probes, mirroring real SOC escalation workflows. That project taught me what I think production AI actually is: mostly engineering --- governance, fallback, observability --- with the modelling sitting on top.

Alongside it, I've deployed a hotel inventory \& procurement platform on the MERN stack --- multi-role (Admin, Manager, Staff, Vendor) with JWT auth, Cloudinary-backed uploads, and Brevo SMTP alerts that pull manual invoice turnaround from hours down to seconds --- and built a fully-offline PDF translation engine using Meta's NLLB-200 accelerated on Apple Silicon (MPS) that brought inference from minutes to seconds. Different domains, but the instinct is the same --- build systems that hold up end-to-end, not just happy-path demos.

Goldman Sachs Engineering is where I'd most like to keep practising that instinct. The combination you describe --- massively scalable systems, low-latency infrastructure, ML alongside financial engineering, and proactive defence against cyber threats --- is exactly the surface area I want to grow on. I'm also drawn to the fact that the engineering org sits at the centre of the business rather than off to the side; the constraints are sharper there, and the feedback loops are real.

I'm currently in my third year of B.E. Information Technology at CBIT Hyderabad (CGPA 9.22), and I'd be available for the 2027 summer internship between my pre-final and final years. My resume is attached, and I'd welcome the chance to talk further.

Thank you for considering my application.

\vspace{12pt}

\noindent Sincerely, \\[10pt]
\textbf{Manuvarthi Seshadri Naidu}

\end{document}
